\section{法的运行}
\begin{enumerate}
  \item 立 | 执、\zd{司}、守 \yjt 法律实施 | 监。
        \begin{enumerate}
          \item 立法。立法法。
          \item 执法。广义 = 司法 + 行政；\zd{狭义}（默认）= 行政。
          \item 司法/法的适用。法院、检察院。\zd{严格的}程序性、合法性。
          \item 守法。积极、消极（观念层面对法律不认同）、作为、不作为。\zd{行使权力 + 履行义务}。
          \item 法律监督。\zd{广义}（默认）、狭义 \yjt 仅指国家机关的监督。监督体系不重要。党的监督 \yjt 社会监督。
        \end{enumerate}
  \item 法学方法论。法律思维培养。总分结构。
        \begin{enumerate}
          \item 法适用的一般原理。总论。
                \begin{enumerate}
                  \item 目标。司法目标。“好”的判决结果。\zd{法理} > 情理。\al{许霆案}。
                        \begin{enumerate}
                          \item 可预测性。依法裁判、形式法治、判决合法。法律效果。“无情”的审判机器。
                          \item 正当性/可接受性。实质法治。公众认同（民意、情理）。社会效果。
                        \end{enumerate}
                  \item 审判。法律发现理论。\zd{个人偏好}。社会因素、心理因素。\zd{论证理论}。\zd{法律和逻辑}。证成 = 说理。
                        \begin{enumerate}
                          \item 大前提 \yjt \zd{法律规范适用、选择}问题，小前提 \yjt 法律事实，结论 \yjt 判决。
                          \item 内部证成。\al{人会死，卡拉是条狗，苏格拉底会死}。法律推理。
                                \begin{itemize}
                                  \item 大小前提正确，\zd{内部逻辑错误}，结论错误。关心逻辑。
                                  \item 内部证成越细致，结论的说服力越强，但不影响结论成立。\al{判决书长短}。
                                \end{itemize}
                          \item 外部证成。\al{人不会死，苏格拉底是人，苏格拉底不会死}。法律解释 + 漏洞填补。
                                \begin{itemize}
                                  \item 大小前提、结论内在逻辑正确，\zd{外部前提错误}，结论错误。关心前提。
                                  \item 必然包含内部证成。搞定前提得符合逻辑。
                                  \item 形式上表现为三段论，\zd{核心是法律解释和法律渊源选择}。
                                \end{itemize}
                        \end{enumerate}
                \end{enumerate}
          \item 法律推理。
                \begin{enumerate}
                  \item 形式推理。形式逻辑。
                        \begin{enumerate}
                          \item \zd{演绎}推理/涵摄（默认）。一般到个别。必然性推理。保真。\al{司法三段论}。
                          \item 归纳推理。个别到一般。或然。不保真。\al{天下乌鸦一般黑；男的没一个好东西}。总结个案。
                          \item 类比推理。个别到个别。或然。\al{判例法国家；指导案例}。对比个案。
                        \end{enumerate}
                  \item 实质推理。
                  \item 设证推理。\al{地上都是水 \yjt 昨天下雨，开豪车 \yjt 家里有矿}。
                        \begin{enumerate}
                          \item 由果推因，或然程度很强，能推翻。
                          \item 高度依赖直觉、经验。\al{立法、侦查、老刑警、苹果落地}。
                        \end{enumerate}
                  \item 反向推理。A \yjt\, R，B \bdy A，B \yjt| R。消极推理。
                        \begin{enumerate}
                          \item 可以有效防止某个规则被随意适用。\al{罪刑法定。法无明文规定不为罪，法无明文不处罚}。
                          \item 明示其一，即否定其余。
                          \item 例外条款。但书规定。无罪辩护。
                        \end{enumerate}
                  \item 当然推理。\zd{理所应当}。定性，轻重 \yjt 再对比 \yjt 结论。
                \end{enumerate}
          \item 法律解释。只有简单、复杂区分。
                \begin{enumerate}
                  \item 对象。法律规定和它的\zd{附随情况}（立法背景、立法资料）。
                  \item 服务于个案。
                  \item 受到\zd{解释学循环}的限制。整体和部分、部分和部分、上下文。\al{跳着看}。\zd{体系解释}。
                  \item 解释方法。自由裁量越来越大。
                        \begin{enumerate}
                          \item 文义解释。字面意思。
                          \item 主观目的解释。\zd{立法者}立法意图。弦外之音、话外之意。\al{你自己看着办吧}。
                          \item 客观目的解释。社会的客观发展需要，新鲜事物，与时俱进。\al{开设赌场，网络赌博}。
                          \item 历史解释。过去的背景、法律、做法。\al{反革命罪}。
                          \item 比较解释。比较法学。\zd{涉外}。
                          \item 体系解释。表述为\zd{联系其他 xx、总则、总论}。考试频率最高。
                        \end{enumerate}
                  \item 选择、位阶问题。
                        \begin{enumerate}
                          \item 单一模式。文义。
                          \item 冲突模式。参考排序。文义 > 体系 > 主观目的 > 历史 > 比较 > 客观目的。
                        \end{enumerate}
                  \item 正式/法定解释 \yjt 立法、司法、行政解释。非正式/任意解释 \yjt 学理、论理解释。
                        \begin{enumerate}
                          \item 立法解释。法律本身、明确界限、补充规定。变相立法。\zd{全人常}，效力 \ydy 法律。
                          \item 行政解释。省法规\zd{应用问题}省政府主管部门，本身找省人常。\zd{国务院及其主管部门}。
                          \item 司法解释。公布 + 要备案。两央两高监省常提出审查申请。\zd{最高法最高检}。
                        \end{enumerate}
                \end{enumerate}
          \item 法律漏洞。\zd{应当规定没规定}。违背立法目的 + 联系整个法律体系。bug。
                \begin{enumerate}
                  \item 分类。
                        \begin{enumerate}
                          \item 全部、部分。漏洞大小。
                          \item 自始 \yjt 立法就有 \yi 明知（有意沉默）\er 不明知（水平问题）。嗣后 \yjt 后来产生。
                          \item 明显 \yjt 说少了，目的扩张，\zd{补充规则}；隐藏 \yjt 说过了，目的限缩，\zd{创设例外}，增加但书。
                        \end{enumerate}
                  \item 填补方法。目的论扩张/限缩（法官造法）。法律规定的范围小于\zd{立法目的}的\zd{潜在/过度包含}。
                        \begin{enumerate}
                          \item \bdy 扩大解释、缩小解释。
                          \item \al{为了顾客的用餐感受，禁止狗进入餐厅，蛇，导盲犬}。
                        \end{enumerate}
                \end{enumerate}
        \end{enumerate}
\end{enumerate}
